%% ============================================================
%%  SECTION 1 — GIỚI THIỆU (Introduction)
%%  Trả lời: Tại sao lại chọn đề tài này? (Motivation)
%% ============================================================
\section{Giới thiệu}
\label{sec:intro}

\subsection{Động lực và tính cấp thiết}

Thị trường cho thuê nhà trọ và căn hộ tại Việt Nam đang phát triển mạnh,
đặc biệt tại các đô thị lớn như Hà Nội và TP. Hồ Chí Minh.
Tuy nhiên, thông tin giá thuê hiện nằm rải rác trên hàng chục trang web với định dạng
và đơn vị không thống nhất; không ít tin đăng chỉ ghi \textit{``giá thỏa thuận''},
khiến người thuê khó so sánh và đánh giá mức giá hợp lý.

Vấn đề này tạo ra ba hướng giá trị thực tiễn mà đồ án hướng tới:
\begin{enumerate}[leftmargin=*, itemsep=3pt]
  \item \textbf{Đối với người thuê:} có thể ước lượng mức giá tham khảo hợp lý
        trước khi đàm phán, tránh bị định giá quá cao.
  \item \textbf{Đối với chủ nhà:} hiểu được mặt bằng khu vực để định giá cạnh tranh.
  \item \textbf{Đối với nghiên cứu:} cung cấp bộ dữ liệu tổng hợp đa nguồn đầu tiên
        về thị trường cho thuê Việt Nam, cùng quy trình tái tạo công khai.
\end{enumerate}

\subsection{Ý nghĩa học thuật}

Từ góc độ Khoa học Dữ liệu, đề tài này bao trùm toàn bộ vòng đời
\textit{Knowledge Discovery in Databases} (KDD): thu thập $\to$ tích hợp $\to$ làm sạch
$\to$ khám phá (EDA) $\to$ mô hình hóa $\to$ triển khai, đặt trọng tâm vào
\textbf{chất lượng dữ liệu} và \textbf{phân tích nhân quả} thay vì chỉ tối ưu chỉ số mô hình.
Hai phát hiện điển hình là \textbf{Nghịch lý Simpson} do cấu trúc nguồn gây ra
(đòi hỏi phân tích phân tầng trước khi kết luận) và việc \textbf{kiểm toán nguồn gốc
dữ liệu} (data provenance) phát hiện 7/12 cờ tiện ích bị \textit{chế tác} ở bước tích hợp —
dẫn tới cơ chế xử lý thiếu-không-ngẫu-nhiên (MNAR) và ràng buộc đơn điệu
cho mô hình triển khai (Mục~\ref{subsec:audit}).

\subsection{Phạm vi bài toán}

\begin{itemize}[leftmargin=*, itemsep=3pt]
  \item \textbf{Input:} diện tích, số phòng ngủ/tắm, thành phố/quận, 12 cờ tiện ích.
  \item \textbf{Output:} giá thuê dự báo (số thực, triệu VNĐ/tháng)
        $\Rightarrow$ \textbf{bài toán Hồi quy}.
  \item \textbf{Phương pháp đánh giá:} MAE, RMSE, R\textsuperscript{2}
        (xem Mục~\ref{sec:eval} để biết cách thiết lập thực nghiệm).
  \item \textbf{Dữ liệu:} 133.586 bản ghi từ 10 nguồn, sau tích hợp và lọc ngoại lệ
        (Mục~\ref{sec:data}).
\end{itemize}

\subsection{Đóng góp chính}

\begin{enumerate}[leftmargin=*, itemsep=3pt]
  \item Xây dựng pipeline thu thập và tích hợp dữ liệu cho thuê đa nguồn
        (10 trang web + bộ dữ liệu công khai).
  \item Phát hiện và phân tích \textbf{Nghịch lý Simpson} do thiên lệch nguồn,
        dẫn tới quyết định thiết kế mô hình có căn cứ.
  \item Thực hiện \textbf{3 thực nghiệm đối sánh} (ablation study) kiểm chứng
        các quyết định tiền xử lý: Feature Engineering vs.\ dữ liệu thô,
        biến mục tiêu gốc vs.\ log-transform, và tác động của ngoại lệ.
  \item \textbf{Kiểm toán sau triển khai} (Mục~\ref{subsec:audit}): truy vết hành vi
        phi lý của mô hình về tận dữ liệu gốc, phát hiện \textbf{7/12 cờ tiện ích bị điền
        ngẫu nhiên} ở bước tích hợp; khắc phục bằng cơ chế \textbf{mask MNAR}
        (thiếu-không-ngẫu-nhiên) và \textbf{ràng buộc đơn điệu}, bảo đảm mô hình
        triển khai hợp lý về kinh tế học (thêm tiện ích không làm giảm giá dự báo).
  \item Triển khai ứng dụng web demo (Streamlit) dự báo giá theo khu vực.
\end{enumerate}
